%% ma: L10-B13-0094
%% lop: 10
%% chuong: 5
%% bai: 13
%% dang: 10.13.12
%% loai: TN4
%% muc: NB
%% dap_an: "C"
%% nguon: "chuyen-de-cac-so-dac-trung-cua-mau-so-lieu-khong-ghep-nhom-toan-10-chuong-trinh-moi.pdf (GV. Phan Nhật Linh) (P03), Bài 02 Dạng 2 Phần I Câu 21"
%% kien_thuc: "Lớp 10 Bài 13 (mốt)"
%% trang_thai: chinh-thuc
%% ghi_chu: "Tự duyệt theo yêu cầu giáo viên 22/10/2026."
\begin{ex}
Cho số liệu thống kê về sản lượng chè thu được trong $1$ năm (kg/sào) của $20$ hộ gia đình:
\begin{center}
\begin{tabular}{|c|c|c|c|c|c|c|c|c|c|}\hline
$111$ & $112$ & $112$ & $113$ & $114$ & $114$ & $115$ & $114$ & $115$ & $116$\\\hline
$112$ & $113$ & $113$ & $114$ & $115$ & $114$ & $116$ & $117$ & $114$ & $115$\\\hline
\end{tabular}
\end{center}
Mốt của mẫu số liệu là
\choice{$M_o=111$}{$M_o=113$}{\True $M_o=114$}{$M_o=117$}
\loigiai{Giá trị $114$ xuất hiện $6$ lần, nhiều nhất. Vậy $M_o=114$. Chọn C.}
\end{ex}
